\textbf{算术模型}是算术语言~$\Lang{L_A}$ 的一个结构。
其中有一个特定的模型，即\textbf{标准模型}~$\Struct{N}$，其论域 $\Domain{N} = \Nat$，
并且对 $\Obj 0$、$\prime$、$+$、$\times$ 和 $<$ 的解释分别由 $0$、$\Nat$ 上的后继函数、加法函数、乘法函数以及小于关系给出。
$\Struct{N}$ 是理论 $\Th{Q}$ 和~$\Th{PA}$ 的模型。

更一般地，$\Lang{L_A}$ 的一个结构当且仅当同构于~$\Struct{N}$ 时称为\textbf{标准的}。
如果两个结构之间存在一个\textbf{同构}，即一个保持常项符号、函数符号和谓词符号解释的双射函数，则它们是同构的。
根据\textbf{同构定理}，同构的结构是\textbf{初等等价的}，也就是说，相同的语句在它们中都为真。
在标准模型中，论域恰好就是所有数字符~$\num{n}$ 的值构成的集合。

与~$\Struct{N}$ 不同构的 $\Th{Q}$ 和 $\Th{PA}$ 的模型被称为\textbf{非标准的}。
在非标准模型中，论域不能被数字符的值穷尽。
$\Struct{M}$ 中的一个元素 $x \in \Domain{M}$，若对所有~$n \in \Nat$ 都有 $x \neq \Value{\num{n}}{M}$，则称为~$\Struct{M}$ 的一个\textbf{非标准元素}。
如果 $\Sat{M}{\Th{Q}}$，非标准元素也必须服从~$\Th{Q}$ 的公理，例如，它们有唯一的后继，可以进行加法和乘法运算，并且可以用~$<$ 来比较。
$\Struct{M}$ 的所有标准元素按~$\Assign{<}{M}$ 都小于所有非标准元素。
非标准模型的存在性由紧致性定理保证，并且对于 $\Th{Q}$，它们相对容易显式地构造出来。
这样的模型可用来表明，比如 $\Th{Q}$ 不足以证明某些语句，例如 $\Th{Q} \Proves/ \lforall[x][\lforall[y][\eq[(x+y)][(y+x)]]]$。
这是通过定义一个非标准模型 $\Struct{M}$ 来实现的，在该模型中，非标准元素不遵守交换律。

$\Th{PA}$ 的非标准模型则不易如此显式地给出。
通过表明 $\Th{PA}$ 可证某些语句，我们可以研究 $\Th{PA}$ 非标准模型中非标准部分的结构。
如果 $\Th{PA}$ 的一个非标准模型~$\Struct{M}$ 是可数的，那么每个非标准元素都属于一个非标准元素的``块''，这些块中的元素在~$\Assign{<}{M}$ 下的序类似于~$\Int$。
而这些块本身的排列方式则类似于~$\Rat$，也就是说，没有最小或最大的块，并且任意两个块之间总存在另一个块。

任何可数模型都同构于一个以~$\Nat$ 为论域的模型。
如果这样一个模型中 $\prime$、$+$、$\times$ 和 $<$ 的解释是可计算函数，我们就称它为一个\textbf{可计算模型}。
标准模型~$\Struct{N}$ 是可计算的，因为 $\Nat$ 上的后继函数、加法函数、乘法函数以及小于关系都是可计算的。
定义 $\Th{Q}$ 的可计算非标准模型是可能的，但 $\Struct{N}$ 是 $\Th{PA}$ 唯一的可计算模型。这就是\textbf{Tannenbaum（坦纳鲍姆）定理}。